% !TeX root = ../../Book2.tex
% 纲要标题：## 第20章　Thom 同构、Euler 类与 Gysin 序列
% 纲要位置：几何学计划纲要.md，第 1769 行。
% 正文状态：初稿；尚未正式审核。

% 纲要内容开始
% 
% > 本章目标与证明范围：从局部纤维定向构造 Thom 类，在相对上同调中证明同构。随后以零截面与球丛解释 Euler 类和 Gysin，不把这些对象分别当作独立定义表。
% 
% 纲要内容结束

\chapter{Thom 同构、Euler 类与 Gysin 序列}
\label{g2:ch:20}
\BookChapterNumber{g2:ch:20}

\begin{chapterabstract}
	纤维的局部方向可以组成总空间上的相对上同调类。本章从盘球丛、Thom 空间及定向局部系统出发，证明任意 CW 基空间上带局部系数的 Thom 同构，说明有限覆盖粘合与无限胞腔的次数控制如何配合。零截面拉回得到 Euler 类，横截截面的零点流形及孤立零点指标给出它的几何解释，并与前文的球丛传递规范比较。最后由盘球丛对的长正合列推导 Gysin 序列，计算 Hopf 丛和圆丛，再用微分形式建立沿纤维积分及其符号约定。
\end{chapterabstract}

\begin{learninggoals}
	\begin{itemize}
		\item 比较盘球丛对、穿孔总空间对与 Thom 商空间，辨认非紧基底及扭系数时各模型的适用范围。
		\item 用纤维的局部同调描述方向输运，解释模二系数和定向局部系统各自如何处理方向反转。
		\item 完整证明局部乘积与整体 Thom 同构，说明无限 CW 基空间为何由有限个维数层控制固定次数。
		\item 从零截面构造 Euler 类，证明它与首障碍、横截零点类及局部指标和的关系。
		\item 用一个具体三维丛辨认 Euler 类消失与非零截面存在之间的区别。
		\item 从长正合列推导 Gysin 序列，计算球丛上同调，并将其积分映射与微分形式的纤维积分对应。
	\end{itemize}
\end{learninggoals}

球面上的向量场为什么不能总是避开零点？第12章把问题放在逐胞腔延拓中，用边界球面的映射度记录首个障碍；第16—17章则把孤立零点附近的映射度相加，得到 Euler 示性数。两种计算都出现同一种局部整数：绕着一个零点走一圈，向量方向绕纤维球面多少次。本章将这个整数放进一张向量丛的相对上同调，从而把延拓、零点和整体上同调联系起来。

\ProseParagraphBreak

先看平凡丛 \(B\times\mathbb R^r\)。每根纤维都带有圆盘对 \((D^r,S^{r-1})\) 的正生成元 \(u_r\)，于是底空间上的类 \(a\) 可以变成相对类 \(a\times u_r\)，次数升高 \(r\)。若丛有扭转，局部正生成元可能在交叠处变号；选择全局方向，或让系数也按同一符号运输，便能把它们拼成 Thom 类。Thom 同构断言，这个升次数构造已经给出盘球丛对的全部上同调。

\ProseParagraphBreak

随后只需对同一个 Thom 类作两种操作。沿零截面拉回得到 Euler 类，它记录一个截面的零点；把相对上同调放回盘球丛对的长正合列，得到 Gysin\footnote{Werner Gysin（1915--1998），瑞士数学家，研究流形纤维化的同调，提出后来以其名字命名的映射与正合列。}序列，其中一支箭头正是乘 Euler 类。这样，三项构造共享同一个局部生成元，而每个映射都能追溯到一个具体的空间对或截面。

\ProseParagraphBreak

本章主要对象是 CW 复形上的有限秩实向量丛；光滑零点与纤维积分部分另说明紧性、横截性及定向条件。第5章的局部系数、第6章的对偶与帽积、第7章的 de Rham 比较、第12章的首障碍及第19章的丛同伦不变性会在用到时逐项回引。Thom 构造与平方运算的联系可参阅\cite[第23章第5--6节]{May1999}，Euler 类与 Thom 同构的经典组织可参阅\cite[第9节，Oriented Bundles and the Euler Class；第10节，The Thom Isomorphism Theorem]{MilnorStasheff1974}。第21章将继续比较由 Thom 平方产生的类与其他示性类构造。

% !TeX root = ../../Book2.tex
% 纲要标题：### 20.1 Thom 空间及定向
% 纲要位置：几何学计划纲要.md，第 1773 行。
% 正文状态：初稿；尚未正式审核。

% 纲要内容开始
% 
% * 盘丛、球丛与相对模型
% * 局部纤维生成元
% * 系数环及扭定向系统
% 
% 纲要内容结束

\section{Thom 空间及定向}
\label{g2:sec:2001}

一根实 \(r\) 维纤维本身可缩，它的普通上同调看不出维数；把零点以外的部分作为相对项后，便出现一个 \(r\) 维生成元。本章把这个局部生成元沿整张向量丛组织起来。首先需要说明采用哪一个相对空间，以及纤维坐标反转方向时生成元如何变化。

\subsection{盘丛、球丛与相对模型}
\label{g2:subsec:2001:01}
本章设 \(B\) 为 CW 复形，\(p:E\to B\) 为秩 \(r\ge1\) 的实向量丛，零截面记为 \(s_0\)。\cref{g2:lem:cw-paracompact-partitions}保证 \(B\) 仿紧；将局部 Euclid 内积按从属单位分解求和，就得到连续纤维内积。这是\secref{g2:sec:1902}构造丛嵌入时已经使用的方法，局部有限性使它同样适用于无限覆盖。以下固定这样一个内积，并记 \(\|v\|\) 为相应范数。

\begin{definition}\label{g2:def:disk-sphere-thom}
	向量丛 \(E\) 的\term[pancong]{盘丛}[disk bundle]、\term[qiucong]{球丛}[sphere bundle]分别为
	\[
		D(E)=\{v\in E:\|v\|\le1\},\qquad S(E)=\{v\in E:\|v\|=1\}.
	\]
	把整个 \(S(E)\) 压成同一个基点 \(*\)，所得带基点空间
	\[
		\operatorname{Th}(E)=D(E)/S(E)
	\]
	称为 \(E\) 的\term[thom-kongjian]{Thom 空间}[Thom space]。另记 \(E^\times=E\setminus s_0(B)\)。
\end{definition}

在局部正交标架中，盘丛和球丛分别是 \(U\times D^r\) 与 \(U\times S^{r-1}\)。Thom 商把所有这些纤维的边界共同送到 \(*\)；各纤维内部的点仍然保留。所以它同时记住纤维的维数及纤维之间的粘合方式。对基底只有一点的情形，这一构造就是 \(D^r/S^{r-1}\cong S^r\)。

\begin{proposition}\label{g2:prop:thom-relative-models}
	对上述向量丛，\((D(E),S(E))\) 是 NDR 对。包含映射
	\[
		(D(E),S(E))\longrightarrow(E,E^\times)
	\]
	在同调和上同调中诱导同构，系数可以取从 \(B\) 拉回的局部系统。对常系数环 \(R\)，商映射还诱导自然同构
	\begin{equation}\label{g2:eq:thom-relative-quotient}
		\widetilde H^k(\operatorname{Th}(E);R)\xrightarrow{\ \cong\ }H^k(D(E),S(E);R).
	\end{equation}
	约化同调也有相同的商空间比较。
\end{proposition}
\begin{proof}
	先给出控制球面附近的径向同伦。对 \(v\in D(E)\)，令 \(\rho=\|v\|\)，并置
	\[
		u(v)=\min\{1,2(1-\rho)\},\qquad
		H(v,t)=\left(1-t+\frac{t}{\max\{\rho,1/2\}}\right)v.
	\]
	分母处处为正，故包括零截面在内该式连续。其范数始终不超过一；当 \(\rho=1\) 时同伦固定 \(v\)，当 \(u(v)<1\) 时有 \(\rho>1/2\)，从而 \(\|H(v,1)\|=1\)。又 \(u^{-1}(0)=S(E)\)。这逐项满足\cref{g2:def:NDR-pair}的条件，所以球丛包含是\cref{g2:thm:NDR-cofibration}所述的闭余纤维化。

	径向映射 \(v\mapsto v/\max\{1,\|v\|\}\) 给出 \(E\) 到 \(D(E)\) 的强形变收缩；线性插值保持每根纤维，且固定盘丛。同样，\(v\mapsto v/\|v\|\) 给出 \(E^\times\) 到 \(S(E)\) 的强形变收缩。因此包含映射在两个空间和两个相对子空间的同调、上同调上分别诱导同构。把这些映射放入对的长正合列，五引理给出所需的相对同构。所有同伦均保持 \(p(v)\) 不变，故拉回系数的识别也是固定的；局部系数情形使用\cref{g2:thm:twisted-homotopy-invariance}的同伦不变性。

	最后核对商空间比较。取球面附近的开领圈
	\[
		N=\{v\in D(E):\|v\|>1/2\}.
	\]
	上面的同伦限制为 \(N\) 到 \(S(E)\) 的强形变收缩，且在商上把 \(N/S(E)\) 收缩到 \(*\)。诱导同伦连续，是因为商映射与局部紧区间 \(I\) 的乘积仍为商映射，见\cref{g2:lem:quotient-interval}。于是长正合列给出
	\[
		\begin{aligned}
			H^k(D(E),N;R)&\xrightarrow{\ \cong\ }H^k(D(E),S(E);R),\\
			H^k(\operatorname{Th}(E),N/S(E);R)&\cong\widetilde H^k(\operatorname{Th}(E);R).
		\end{aligned}
	\]
	商映射在 \(D(E)\setminus S(E)\) 上是到 \(\operatorname{Th}(E)\setminus\{*\}\) 的同胚。分别从上述两个对中切除 \(S(E)\) 和 \(\{*\}\)，所得对因而同胚。切除条件成立：这两个闭集分别包含于开集 \(N\) 和 \(N/S(E)\) 内。由\cref{g2:prop:relative-cohomology-exactness}的上同调切除，两个对的上同调相同。合并这些同构便得到式\eqref{g2:eq:thom-relative-quotient}，而且合成映射正是商映射诱导的映射。同调的证明逐步相同，所用切除是\cref{g2:thm:singular-excision}。
\end{proof}

盘球丛模型适合使用长正合列，\((E,E^\times)\) 适合讨论零截面附近的几何，而 Thom 空间把相对问题写成一个带基点空间的约化上同调。以上论证说明这些模型为何给出同一份常系数信息。局部系数将在相对模型上保留：从 \(B\) 拉回的局部系统一般不能在压缩整个 \(S(E)\) 后成为商空间上的局部系统。

\ProseParagraphBreak

内积的选择不会改变 Thom 空间的同胚类型。若两组范数为 \(\|\cdot\|_0,\|\cdot\|_1\)，则
\[
	v\longmapsto
	\begin{cases}
		\dfrac{\|v\|_0}{\|v\|_1}v,&v\ne0,\\
		0,&v=0
	\end{cases}
\]
把第一组盘球丛送到第二组盘球丛，交换下标便得逆。在任何局部标架中，两组正定内积在每个基点附近具有统一的范数比较常数，所以该映射在零截面也连续。它因此诱导 Thom 空间的同胚。若 \(B\) 紧，则 \(D(E)\) 紧，且 \(v\mapsto v/(1-\|v\|)\) 将开盘丛同胚到 \(E\)，这时 \(\operatorname{Th}(E)\) 正是总空间 \(E\) 的一点紧化。非紧基底时使用上述盘球丛商定义；\cref{g2:ex:noncompact-thom-not-compactification}将具体比较两种构造。

\subsection{局部纤维生成元}
\label{g2:subsec:2001:02}
固定 \(b\in B\)，记 \(D_b=D(E)\cap E_b\)、\(S_b=S(E)\cap E_b\)。相对球的计算给出
\begin{equation}\label{g2:eq:fiber-orientation-group}
	H_j(D_b,S_b;\mathbb Z)\cong
	\begin{cases}
		\mathbb Z,&j=r,\\
		0,&j\ne r.
	\end{cases}
\end{equation}
这也可由 \(D_b/S_b\cong S^r\) 及\cref{g2:thm:sphere-homology}得到。选一组有序线性坐标，就选出顶维相对基本类 \(o_b=[D_b,S_b]\)；反转坐标方向时它变成 \(-o_b\)。由上一小节的相对比较，它也确定 \(H_r(E_b,E_b\setminus\{0\};\mathbb Z)\) 的一个生成元。

\ProseParagraphBreak

对于任意 Abel 群 \(A\)，上同调万有系数定理给出自然识别
\begin{equation}\label{g2:eq:fiber-top-cohomology-hom}
	H^r(D_b,S_b;A)\cong\operatorname{Hom}_{\mathbb Z}\bigl(H_r(D_b,S_b;\mathbb Z),A\bigr),
\end{equation}
其余次数的相对上同调为零。这里顶维相对同调是自由群，且其余相对同调消失，所以万有系数定理中的 \(\operatorname{Ext}\) 项为零，见式\eqref{g2:eq:uct-cohomology}。当 \(A=\mathbb Z\) 时，与 \(o_b\) 对偶的生成元 \(u_b\) 由
\[
	\langle u_b,o_b\rangle=1
\]
刻画。方向改变使 \(o_b\) 和 \(u_b\) 同时变号。这一配对将是把局部方向转成 Thom 类的标准化条件。

\ProseParagraphBreak

球面取向采用外法向在前的边界约定。长正合列的边界映射于是满足 \(\partial o_b=[S_b]\)。秩一时这个公式读成
\[
	\partial[-1,1]=[1]-[-1]\in\widetilde H_0(S^0;\mathbb Z).
\]
若球面的约化顶维上同调类 \(\omega_b\) 在 \([S_b]\) 上取值一，则对的连接映射满足 \(\delta\omega_b=u_b\)：在链上有 \(\langle\delta\omega_b,o_b\rangle=\langle\omega_b,\partial o_b\rangle=1\)。这个约定固定了后面球丛长正合列中纤维积分的正负号。

\begin{definition}\label{g2:def:bundle-orientation-system}
	向量丛 \(E\) 的\term[xiangliangcong-dingxiang-jubuxitong]{定向局部系统}[orientation local system]记为 \(\mathcal O_E\)，其在 \(b\) 处的群为
	\[
		\mathcal O_{E,b}=H_r(E_b,E_b\setminus\{0\};\mathbb Z).
	\]
	局部平凡化给出这些群之间的识别，沿路径依次使用局部识别便得到输运。若可为每个 \(b\) 选取生成元 \(o_b\in\mathcal O_{E,b}\)，使它们在局部平凡化中相容，则称 \(E\) \term[xiangliangcong-kedingxiang]{可定向}[orientable]；这样一组选择称为 \(E\) 的一个定向。
\end{definition}

这一定向描述每根向量空间纤维的方向，基底本身不需要是流形，也不需要已有定向。它与\chapref{g2:ch:05}中流形的定向局部系统使用同一种局部同调思想；此处取的是纤维在零点的局部同调。

\begin{proposition}\label{g2:prop:bundle-orientation-transport}
	若两组局部坐标的过渡函数为 \(g_{\alpha\beta}:U_\alpha\cap U_\beta\to\operatorname{GL}_r(\mathbb R)\)，则 \(\mathcal O_E\) 的相应过渡函数为
	\[
		\operatorname{sgn}\det g_{\alpha\beta}\in\{1,-1\}.
	\]
	这些函数定义秩一的整数局部系统，其路径输运只依赖端点固定的路径同伦类。在每个路径连通分支上，\(E\) 可定向当且仅当所有闭路在 \(\mathcal O_E\) 上的输运都是恒等。
\end{proposition}
\begin{proof}
	先考察一个可逆线性映射 \(A:\mathbb R^r\to\mathbb R^r\)。极分解把它写成正交矩阵与正定矩阵的乘积；正定矩阵可沿 \((1-t)P+tI\) 连到恒等。行列式为正的正交矩阵由二维平面旋转组成，可以把各旋转角连续缩到零。因此 \(\det A>0\) 时，\(A\) 与恒等映射通过可逆线性映射同伦，在相对顶维同调上作用为 \(+1\)。若 \(\det A<0\)，先复合一次坐标反射便化为前一种情形；反射在相对基本类上作用为 \(-1\)，与\cref{g2:prop:reflection-antipodal-degree}的球面度计算一致。故所求作用正是 \(\operatorname{sgn}\det A\)。

	行列式连续且不为零，所以其符号在交集上局部常值。过渡函数的乘法关系经行列式符号仍成立，因而确实得到局部系统。用这些局部常值识别沿路径输运，细分路径或改变局部平凡化均由过渡关系保证不改变合成；对路径同伦，把参数方形细分到每个小方形落入一个平凡化域，小方形内的输运相消，得到端点固定的同伦不变性。这也就是\secref{g2:sec:0501}中局部系统的路径描述。

	若已有相容生成元，任何闭路都把它送回自身，故输运恒等。反过来，在一个分支中固定 \(b_0\) 及一个生成元 \(o_{b_0}\)，沿从 \(b_0\) 到 \(b\) 的路径输运它。两条路径之差是一条闭路，所以假设保证所得 \(o_b\) 与路径无关；局部平凡化中的输运恒定，又保证这些生成元相容。CW 复形局部道路连通，因此上述逐分支的构造给出所需定向。
\end{proof}

\begin{example}\label{g2:exm:complex-bundle-real-orientation}
	每张复秩 \(m\) 向量丛的底层实秩 \(2m\) 丛都有自然定向。对一组复基 \(e_1,\ldots,e_m\)，取实有序基
	\[
		(e_1,\mathrm i e_1,e_2,\mathrm i e_2,\ldots,e_m,\mathrm i e_m).
	\]
	若复坐标变换矩阵为 \(A\)，其底层实矩阵的行列式为 \(\det_{\mathbb R}A=|\det_{\mathbb C}A|^2>0\)。这个公式可在复化后看出：实矩阵的复化分解为 \(A\) 与 \(\overline A\) 两块，行列式为 \(\det A\cdot\overline{\det A}\)。因此更换复标架不会改变上述实方向，局部方向自动粘合。
\end{example}

\subsection{系数环及扭定向系统}
\label{g2:subsec:2001:03}
取交换有幺环 \(R\)，并令
\[
	\mathcal O_{E,R}=\mathcal O_E\otimes_{\mathbb Z}R.
\]
这是秩一自由 \(R\) 模的局部系统，过渡函数仍为 \(\pm1\)，现将它们视作 \(R\) 中的单位。一个 \(R\) 定向是 \(\mathcal O_{E,R}\) 与常局部系统 \(\underline R\) 的同构，等价地，是一组相容的纤维 \(R\) 模基元。整数定向经系数变换给出 \(R\) 定向；若 \(2=0\) 于 \(R\)，则两个整数方向都给出同一个模 \(R\) 基元，故每张实向量丛都有规范的 \(R\) 定向。特别地，模二系数完全消除了方向反转造成的符号。

\begin{example}\label{g2:exm:mobius-orientation-system}
	设 \(L\to S^1\) 为 Möbius 线丛，其矩形模型沿两端按 \((0,t)\sim(1,-t)\) 粘合。沿基圆一周，纤维的相对基本类变成其相反数，所以 \(\mathcal O_L\) 的单值作用为 \(-1\)，它没有整数定向。若一个类 \(u\in H^1(D(L),S(L);\mathbb Z)\) 在每根纤维上都是生成元，则沿基圆输运后其限制必须仍为自身，同时又要变号。这要求一个整数生成元等于其相反数，矛盾。用 \(\mathbb F_2\) 系数时两个符号相同，方向障碍消失。\secref{g2:sec:1304}已经利用这一模二 Thom 类研究了平方运算，见\cref{g2:exm:mobius-thom-square}。
\end{example}

为了对所有实丛统一处理整数系数，可以让系数也随纤维方向一起变化。固定 \(b\)，式\eqref{g2:eq:fiber-top-cohomology-hom}给出规范识别
\begin{equation}\label{g2:eq:twisted-fiber-endomorphism}
	H^r(D_b,S_b;\mathcal O_{E,R,b})
	\cong\operatorname{Hom}_{\mathbb Z}(\mathcal O_{E,b},\mathcal O_{E,R,b})
	\cong\operatorname{End}_R(\mathcal O_{E,R,b}).
\end{equation}
第二个同构把群同态唯一延拓为 \(R\) 线性映射。右端有一个不依赖任何基元的元素：恒等映射。整数系数时它就是 \(\operatorname{id}_{\mathcal O_{E,b}}\)；一般系数时，在中间一项对应 \(m\mapsto m\otimes1\)。

\begin{definition}\label{g2:def:twisted-thom-class}
	类
	\[
		U_E\in H^r(D(E),S(E);p^*\mathcal O_{E,R})
	\]
	若在每个纤维上的限制都对应式\eqref{g2:eq:twisted-fiber-endomorphism}中的恒等映射，就称为 \(E\) 的\term[niu-thom-lei]{扭 Thom 类}[twisted Thom class]。需要区分系数环时记为 \(U_{E,R}\)。若选定 \(E\) 的 \(R\) 定向 \(\mathcal O_{E,R}\cong\underline R\)，由它得到的常系数类记为 \(u_E\)，称为这个定向的 Thom 类；其纤维限制在选定的相对基本类上取值一。
\end{definition}

这一标准化把两次变号配成了方向无关的条件。具体地，选一个整数方向 \(o_b\)，并令 \(u_b(o_b)=1\)，则局部扭类写成 \(u_b\otimes(o_b\otimes1)\)。更换方向时两个因子都变号，张量保持不变；在式\eqref{g2:eq:twisted-fiber-endomorphism}中，这正是同一个恒等映射。在 Möbius 例中，它也沿闭路保持不变。局部条件能否确定一个全局上同调类，需要进一步的粘合证明；下一节的\cref{g2:thm:thom-isomorphism}将证明这样的类存在、唯一，并说明乘以它为何给出次数平移。

\begin{proposition}\label{g2:prop:bundle-orientation-square}
	存在规范的局部系统同构
	\begin{equation}\label{g2:eq:bundle-orientation-square}
		\mathcal O_{E,R}\otimes_R\mathcal O_{E,R}\xrightarrow{\ \cong\ }\underline R.
	\end{equation}
	在任一整数方向 \(o_b\) 下，它将 \((o_b\otimes a)\otimes(o_b\otimes c)\) 送到 \(ac\)。
\end{proposition}
\begin{proof}
	每个纤维都是秩一自由 \(R\) 模，所以该规则给出一个同构。把 \(o_b\) 换成 \(-o_b\) 时，同一元素的两个系数都变成相反数，乘积不变，故映射与方向选择无关。沿路径输运时，两个定向因子具有同一符号，其乘积为 \((\pm1)^2=1\)，因而这些纤维同构与输运相容，给出局部系统同构。
\end{proof}

这与\cref{g2:prop:orientation-square}中流形定向系统的平方配对是同一种构造。它解释了后面公式中为什么会出现一个或两个定向因子：乘一个扭 Thom 类会添上 \(\mathcal O_{E,R}\)，两个相同的定向因子则可按式\eqref{g2:eq:bundle-orientation-square}规范消去。此后的杯积按“底空间类在前、Thom 类在后”排列，记作 \(p^*a\smile U_E\)。这一顺序同时确定次数交换可能带来的符号。
% 纲要条目排为编号小节标题，并进入目录与书签；说明段仍保留于纲要注释中。

% !TeX root = ../../Book2.tex
% 纲要标题：### 20.2 Thom 同构的证明
% 纲要位置：几何学计划纲要.md，第 1729 行。
% 正文状态：未撰写；本文件目前是框架。

% 纲要内容开始
% 
% * 局部乘积上的杯积同构
% * Mayer–Vietoris 粘合
% * 一般基空间的有限维或穷竭控制
% * 自然性与直和相容
% 
% 纲要内容结束

\section{Thom 同构的证明}
\label{g2:sec:2002}

待撰写

\subsection{局部乘积上的杯积同构}
\label{g2:subsec:2002:01}
待撰写

\subsection{Mayer–Vietoris 粘合}
\label{g2:subsec:2002:02}
待撰写

\subsection{一般基空间的有限维或穷竭控制}
\label{g2:subsec:2002:03}
待撰写

\subsection{自然性与直和相容}
\label{g2:subsec:2002:04}
待撰写
% 纲要条目排为编号小节标题，并进入目录与书签；说明段仍保留于纲要注释中。

% !TeX root = ../../Book2.tex
% 纲要标题：### 20.3 Euler 类与零点
% 纲要位置：几何学计划纲要.md，第 1736 行。
% 正文状态：未撰写；本文件目前是框架。

% 纲要内容开始
% 
% * 零截面拉回
% * 横截截面的零点类
% * 局部指标与 Euler 数
% * 零 Euler 类不足以消除所有高阶障碍
% 
% 纲要内容结束

\section{Euler 类与零点}
\label{g2:sec:2003}

待撰写

\subsection{零截面拉回}
\label{g2:subsec:2003:01}
待撰写

\subsection{横截截面的零点类}
\label{g2:subsec:2003:02}
待撰写

\subsection{局部指标与 Euler 数}
\label{g2:subsec:2003:03}
待撰写

\subsection{零 Euler 类不足以消除所有高阶障碍}
\label{g2:subsec:2003:04}
待撰写
% 纲要条目排为编号小节标题，并进入目录与书签；说明段仍保留于纲要注释中。

% !TeX root = ../../Book2.tex
% 纲要标题：### 20.4 Gysin 序列及积分沿纤维
% 纲要位置：几何学计划纲要.md，第 1793 行。
% 正文状态：未撰写；本文件目前是框架。

% 纲要内容开始
% 
% * 从盘球丛对的长正合列推导
% * 乘 Euler 类的映射
% * Hopf 丛及球丛计算
% * 定向纤维积分的上同调模型
% 
% ---
% 
% 纲要内容结束

\section{Gysin 序列及积分沿纤维}
\label{g2:sec:2004}

待撰写

\subsection{从盘球丛对的长正合列推导}
\label{g2:subsec:2004:01}
待撰写

\subsection{乘 Euler 类的映射}
\label{g2:subsec:2004:02}
待撰写

\subsection{Hopf 丛及球丛计算}
\label{g2:subsec:2004:03}
待撰写

\subsection{定向纤维积分的上同调模型}
\label{g2:subsec:2004:04}
待撰写
% 纲要条目排为编号小节标题，并进入目录与书签；说明段仍保留于纲要注释中。


\begin{chaptersummary}
	盘球丛对把每根纤维的方向记录为一个最高次相对生成元；定向局部系统同时保存这些生成元沿道路的符号变化。局部乘积给出次数平移，有限覆盖通过 Mayer--Vietoris 粘合，一般 CW 基空间则由相对胞腔复形控制。所得 Thom 类唯一，并与拉回、直和及系数变换相容；任意局部系数的上同调都可由与它作杯积移入盘球丛对。

	\ProseParagraphBreak

	零截面拉回 Thom 类得到 \(e_{\mathrm{Th}}(E)\)。它等于球丛的首障碍，也等于横截零点流形的对偶类；秩与底流形维数相等时，求值就是孤立零点的指标和。按前文已固定的传递符号，它与球丛的 \(e(S(E))\) 相差一个负号。Euler 类消失只消除首障碍，高维胞腔还可能阻止非零截面存在。

	\ProseParagraphBreak

	盘球丛对的长正合列经 Thom 同构改写为 Gysin 序列：忘去相对条件对应乘 Euler 类，球丛上的类经连接映射再作 Thom 逆同构，得到降低纤维维数的积分映射。Hopf 丛与二球面上的圆丛展示了它的计算作用。按底空间微分在前、纤维微分在后的顺序，形式积分满足无附加符号的投影公式，并与规范化后的拓扑积分一致；整数序列还保留实微分形式看不到的扭信息。
\end{chaptersummary}

\BookExercises

\begin{exercise}\label{g2:ex:noncompact-thom-not-compactification}
	取 \(\mathbb R\) 上的平凡实线丛 \(E=\mathbb R\times\mathbb R\)。直接从盘球丛模型证明 \(\operatorname{Th}(E)\simeq S^1\)，并计算其约化整数上同调。再与总空间 \(E\) 的一点紧化比较，说明非紧基底时两种构造确实可能不同。
\end{exercise}
\begin{solution}
	盘丛和球丛分别是 \(\mathbb R\times[-1,1]\) 与 \(\mathbb R\times\{-1,1\}\)。同伦
	\[
		F_s(x,t)=((1-s)x,t),\qquad 0\le s\le1,
	\]
	始终把球丛映入球丛，固定中间纤维 \(\{0\}\times[-1,1]\)，所以下降为 Thom 空间上的同伦。下降后的同伦连续，由\cref{g2:lem:quotient-interval}的积商性质保证。终映射的像为中间纤维的商 \([-1,1]/\{-1,1\}\cong S^1\)，因而得到一个强形变收缩。于是
	\[
		\widetilde H^k(\operatorname{Th}(E);\mathbb Z)\cong
		\begin{cases}
			\mathbb Z,&k=1,\\
			0,&k\ne1.
		\end{cases}
	\]
	总空间则同胚于 \(\mathbb R^2\)，其一点紧化为 \(S^2\)，约化整数上同调的非零项在二次。因此它与 \(\operatorname{Th}(E)\) 连同伦等价都不成立。盘球丛商只将纤维的单位边界共同压成一点；总空间的一点紧化还把沿非紧基底逃向无穷远的方向纳入无穷远点的邻域条件。
\end{solution}

\begin{exercise}\label{g2:ex:doubled-bundle-orientation}
	设 \(E\to B\) 为实秩 \(r\) 向量丛。对任意局部实标架 \(e_1,\ldots,e_r\)，用有序标架
	\[
		((e_1,0),\ldots,(e_r,0),(0,e_1),\ldots,(0,e_r))
	\]
	给 \(E\oplus E\) 定向。证明该定向不依赖局部标架，所以即使 \(E\) 不可定向，\(E\oplus E\) 仍有这一规范定向。再用 \((v,w)\mapsto v+\mathrm i w\) 识别它与复化丛 \(E\otimes_{\mathbb R}\mathbb C\) 的底层实丛，比较题给定向与复结构的自然定向。
\end{exercise}
\begin{solution}
	若原标架的变换矩阵是 \(A\in\operatorname{GL}_r(\mathbb R)\)，则直和中这组有序标架的变换矩阵为 \(\operatorname{diag}(A,A)\)，行列式是 \((\det A)^2>0\)。所以所有局部方向相容，且更换任何局部标架都不改变定向。这一论证不要求 \(\det A\) 本身为正。

	在题给识别下，直和定向对应实有序基
	\[
		(e_1,\ldots,e_r,\mathrm i e_1,\ldots,\mathrm i e_r),
	\]
	而复结构的自然定向对应交错排列 \((e_1,\mathrm i e_1,\ldots,e_r,\mathrm i e_r)\)。把前一个排列变为后一个时，\(\mathrm i e_j\) 需要跨过 \(r-j\) 个实基向量，总交换次数为
	\[
		(r-1)+(r-2)+\cdots+1=\frac{r(r-1)}2.
	\]
	故两个定向的比较符号为 \((-1)^{r(r-1)/2}\)：\(r\equiv0,1\pmod4\) 时相同，\(r\equiv2,3\pmod4\) 时相反。两种定向都与标架无关，但所采用的有序基排列不同。
\end{solution}

\begin{exercise}\label{g2:ex:orientation-over-z-four}
	令 \(L\to S^1\) 为 Möbius 线丛。分别对 \(R=\mathbb Z\)、\(\mathbb Z/4\)、\(\mathbb F_2\)，计算 \(H^0(S^1;\mathcal O_{L,R})\)，并判断 \(L\) 是否 \(R\) 可定向。说明存在非零平行截面为何还不足以给出 \(R\) 定向。
\end{exercise}
\begin{solution}
	这三个定向系统沿正向圆周的单值作用都是乘以 \(-1\)。\cref{g2:prop:circle-local-coefficients}给出圆周局部系数零次上同调为单值作用的不变量，因此
	\[
		H^0(S^1;\mathcal O_{L,R})=\ker(-1-1:R\longrightarrow R)=\{a\in R:2a=0\}.
	\]
	在三种系数中依次得到 \(0\)、\(\{0,2\}\subset\mathbb Z/4\)、\(\mathbb F_2\)。一个 \(R\) 定向要求平行截面的值为秩一自由 \(R\) 模的基元，选定一个局部基后，这等价于该值是 \(R\) 中的单位。

	整数系数没有非零平行截面，所以不可定向。\(\mathbb Z/4\) 系数虽有非零平行截面 \(2\)，但它不是单位，所生成的子模只有 \(\{0,2\}\)，不能成为纤维模的一组基；故仍不可定向。\(\mathbb F_2\) 系数的平行截面 \(1\) 是单位，因而给出规范定向。第二个例子说明，在带有零因子的系数环上，“非零”与“生成整个秩一自由模”是两个不同的条件。
\end{solution}

\begin{exercise}\label{g2:ex:mobius-thom-cohomology}
	设 \(L\to S^1\) 为 Möbius 线丛。分别计算 \(H^k(D(L),S(L);\mathbb Z)\)、\(H^k(D(L),S(L);p^*\mathcal O_L)\) 与 \(H^k(D(L),S(L);\mathbb F_2)\)。再用盘球丛对的长正合列独立验证第一组结果，并解释其二次整数类为什么不能充当 Thom 类。
\end{exercise}
\begin{solution}
	由式\eqref{g2:eq:twisted-thom-isomorphism}及 \(\mathcal O_L\otimes\mathcal O_L\cong\mathbb Z\)，普通整数相对上同调对应 \(H^{k-1}(S^1;\mathcal O_L)\)，扭系数相对上同调对应 \(H^{k-1}(S^1;\mathbb Z)\)。定向系统的圆周胞腔上链为
	\[
		\mathbb Z\xrightarrow{\,-2\,}\mathbb Z,
	\]
	所以它的零次上同调为零、一次为 \(\mathbb Z/2\)。于是
	\[
	\begin{array}{c|ccc}
		&k=1&k=2&\text{其余 }k\\\hline
		H^k(D(L),S(L);\mathbb Z)&0&\mathbb Z/2&0\\
		H^k(D(L),S(L);p^*\mathcal O_L)&\mathbb Z&\mathbb Z&0\\
		H^k(D(L),S(L);\mathbb F_2)&\mathbb F_2&\mathbb F_2&0
	\end{array}
	\]
	最后一行使用模二方向系统平凡及圆周的模二上同调。

	独立验证时，\(D(L)\) 收缩到基圆，\(S(L)\) 也是一个圆，而其投影绕基圆两次。故限制 \(H^1(D(L);\mathbb Z)\to H^1(S(L);\mathbb Z)\) 为乘二；两者的零次限制为恒等。长正合列于是给出
	\[
		0\longrightarrow H^1(D(L),S(L);\mathbb Z)\longrightarrow\mathbb Z
		\xrightarrow{\,2\,}\mathbb Z\longrightarrow H^2(D(L),S(L);\mathbb Z)\longrightarrow0,
	\]
	以及其他次数消失，与上表相符。秩一丛的 Thom 类须是一次类，并在每根区间对上取生成元；这里普通整数一次相对群已经为零。二次的扭元素限制到纤维后为零，因为纤维对的二次上同调为零，所以它不能替代 Thom 类。
\end{solution}

\begin{exercise}\label{g2:ex:infinite-thom-cochains}
	令 \(B=\bigvee_{j\ge1}S^1_j\) 具有 CW 弱拓扑，\(E=B\times\mathbb R^2\) 取标准定向。计算 \(\widetilde H^k(\operatorname{Th}(E);\mathbb Z)\)。找出一个三次类，使它在每个圆周上限制后都非零，并说明用有限子楔和的上同调按补零映射取直极限会漏掉什么。
\end{exercise}
\begin{solution}
	\(B\) 有一个零胞腔及可数多个一胞腔，胞腔边界全为零。因此
	\[
		H^0(B;\mathbb Z)=\mathbb Z,\qquad H^1(B;\mathbb Z)=\prod_{j\ge1}\mathbb Z,
	\]
	更高次为零。Thom 同构与常系数商比较给出
	\[
		\widetilde H^k(\operatorname{Th}(E);\mathbb Z)=
		\begin{cases}
			\mathbb Z,&k=2,\\
			\prod_{j\ge1}\mathbb Z,&k=3,\\
			0,&\text{其他次数}.
		\end{cases}
	\]
	取一次胞腔上链 \(a=(1,1,\ldots)\)。它是上循环且不为余边界，\(\Phi_E(a)\) 在限制到任意 \(S^1_j\) 上的丛时都是三次群的生成元，这由 Thom 自然性可知。

	前 \(N\) 个圆周的上同调为 \(\mathbb Z^N\)。若按题给补零映射取直极限，得到的是 \(\bigoplus_{j\ge1}\mathbb Z\)，只含有限支撑序列；它不含 \((1,1,\ldots)\)。实际由子空间包含诱导的上同调映射是反向的限制，而非补零映射。胞腔计算使用全部胞腔系数的直积，因而保留了这个类，也说明固定次数的有限骨架控制不等于有限胞腔支撑。
\end{solution}

\begin{exercise}\label{g2:ex:thom-space-product-euler}
	设 \(E\to B\) 为已定向实秩 \(r\ge1\) 丛，取系数环 \(R\)，记 \(e=e_{\mathrm{Th}}(E)\)。证明 Thom 同构一般只是 \(H^*(B;R)\) 模的次数平移同构；其乘法满足
	\[
		\Phi_E(a)\smile\Phi_E(b)=(-1)^{r\deg b}\Phi_E(a\smile b\smile e).
	\]
	特别计算 \(u_E^2\)，并说明为什么平凡正秩丛的 Thom 空间约化上同调中任意两个类的乘积都为零。
\end{exercise}
\begin{solution}
	记 \(j^*:H^*(D(E),S(E);R)\to H^*(D(E);R)\) 为忘记相对条件的映射。零截面与投影互为同伦逆，而 Euler 类定义给出 \(s_0^*j^*u_E=e\)，故
	\[
		j^*u_E=p^*e.
	\]
	相对杯积与绝对—相对杯积相容：两个相对上链都在 \(S(E)\) 上为零，将其中一个视为绝对上链，不改变它们的杯积。因此
	\[
		u_E\smile u_E=j^*u_E\smile u_E=p^*e\smile u_E=\Phi_E(e).
	\]
	对一般齐次类，把次数 \(r\) 的 \(u_E\) 越过 \(p^*b\)，得到
	\[
		(p^*a\smile u_E)\smile(p^*b\smile u_E)
		=(-1)^{r\deg b}p^*(a\smile b\smile e)\smile u_E,
	\]
	即题给公式。非齐次类按齐次分量展开即可。

	平凡正秩丛有常值非零截面，由\cref{g2:prop:euler-basic-properties}得 \(e=0\)，所以上式恒为零。经商比较，这就是 Thom 空间约化上同调的乘法；它不意味着底空间本身的上同调乘法为零。Thom 同构保留底空间类作用在相对类上的模结构，约化环的内部乘法则额外记住 Euler 类。
\end{solution}

\begin{exercise}\label{g2:ex:odd-rank-euler-torsion}
	设 \(E\to B\) 是奇数秩的定向实向量丛。证明 \(2e_{\mathrm{Th}}(E)=0\)。据此说明：若 \(H^{\operatorname{rank}E}(B;\mathbb Z)\) 没有非零的二阶扭元，则 Euler 类为零；但这一点仍不保证存在非零截面。再用这个结论与\cref{g2:cor:tangent-euler-characteristic}证明闭定向奇维流形的 Euler 示性数为零。
\end{exercise}
\begin{solution}
	记 \(E^{\mathrm{op}}\) 为反转纤维定向后的同一个丛。逐纤维映射 \(-I:E\to E^{\mathrm{op}}\) 保持定向：它相对于原定向的行列式符号为 \((-1)^{\operatorname{rank}E}=-1\)，而目标又反转了一次定向。Euler 类对保持定向的丛同构不变，因此
	\[
		e_{\mathrm{Th}}(E)=e_{\mathrm{Th}}(E^{\mathrm{op}})=-e_{\mathrm{Th}}(E).
	\]
	于是 \(2e_{\mathrm{Th}}(E)=0\)。若相应上同调群没有非零二阶扭元，就有 \(e_{\mathrm{Th}}(E)=0\)。\cref{g2:ex:zero-euler-no-section}中的 \(S^4\) 上定向三维丛满足这个条件，却没有非零截面，故消失条件不充分。

	若 \(M\) 是闭定向奇维流形，则整数求值满足 \(2\langle e_{\mathrm{Th}}(\mathrm TM),[M]\rangle=0\)，而整数群没有非零二阶扭元，所以这个求值为零。由\cref{g2:cor:tangent-euler-characteristic}，它恰是 \(\chi(M)\)。
\end{solution}

\begin{exercise}\label{g2:ex:product-tangent-euler}
	设 \(M,N\) 是闭定向光滑流形，维数分别为正整数 \(m,n\)。乘积按 \(M\) 在前、\(N\) 在后定向。用切丛分解和 Euler 类证明
	\[
		e_{\mathrm{Th}}(\mathrm T(M\times N))
		=\pi_M^*e_{\mathrm{Th}}(\mathrm TM)\smile\pi_N^*e_{\mathrm{Th}}(\mathrm TN),
		\qquad\chi(M\times N)=\chi(M)\chi(N).
	\]
	进一步，若 \(X,Y\) 分别是 \(M,N\) 上仅有非退化零点的向量场，计算乘积向量场 \((X,Y)\) 在 \((x,y)\) 处的指标。
\end{exercise}
\begin{solution}
	乘积坐标给出保持所选定向的丛同构
	\[
		\mathrm T(M\times N)\cong\pi_M^*\mathrm TM\oplus\pi_N^*\mathrm TN.
	\]
	由拉回自然性和直和公式，得到所求的第一式。乘积基本类为 \([M]\times[N]\)，而按同一因子次序的外积满足
	\[
		\bigl\langle\pi_M^*a\smile\pi_N^*b,[M]\times[N]\bigr\rangle
		=\langle a,[M]\rangle\langle b,[N]\rangle
	\]
	对 \(a\in H^m(M;\mathbb Z)\)、\(b\in H^n(N;\mathbb Z)\) 成立。这里没有交换两个因子，故没有附加符号。将两个 Euler 类代入，并在三个流形上使用\cref{g2:cor:tangent-euler-characteristic}，便得第二式。

	乘积向量场的零点恰为两个零点集的笛卡儿积。在正向乘积坐标下，其线性化为分块对角矩阵
	\[
		\begin{pmatrix}\mathrm DX_x&0\\0&\mathrm DY_y\end{pmatrix}.
	\]
	因此每个零点非退化，且
	\[
		\operatorname{ind}_{(x,y)}(X,Y)
		=\operatorname{sgn}(\det\mathrm DX_x\det\mathrm DY_y)
		=\operatorname{ind}_x(X)\operatorname{ind}_y(Y).
	\]
	对所有零点求和，得到两个指标和的乘积，也从零点计数直接解释了 Euler 数的乘法公式。
\end{solution}

\begin{exercise}\label{g2:ex:pullback-tangent-sphere-euler}
	设 \(f:S^2\to S^2\) 是次数为 \(d\in\mathbb Z\) 的光滑映射，给 \(E=f^*\mathrm TS^2\) 取拉回定向。计算 \(e_{\mathrm{Th}}(E)\)，判定 \(E\) 何时存在非零截面。若 \(s\) 是与零截面横截的截面，证明它至少有 \(2|d|\) 个零点，并且零点总数为偶数。
\end{exercise}
\begin{solution}
	令 \(\omega\in H^2(S^2;\mathbb Z)\) 在正基本类上取值一。由自然性、映射度的定义和\cref{g2:ex:euler-tangent-spheres-surfaces}，
	\[
		e_{\mathrm{Th}}(E)=f^*(2\omega)=2d\omega.
	\]
	若有非零截面，Euler 类必须为零，故 \(d=0\)。反过来，若 \(d=0\)，Euler 类为零，而底空间维数等于丛的秩；\cref{g2:thm:euler-primary-comparison}与\cref{g2:thm:sphere-primary-euler}说明首障碍消失后即可得到全空间上的单位截面。因此非零截面存在当且仅当 \(d=0\)。这里先得到连续截面；若要求光滑，可将其充分接近地光滑逼近，紧性保证逼近后仍非零。

	设横截零点中指标为 \(+1,-1\) 的个数分别为 \(N_+,N_-\)。\cref{g2:thm:euler-index-evaluation}给出 \(N_+-N_-=2d\)，故
	\[
		N_++N_-\ge|N_+-N_-|=2|d|,
		\qquad N_++N_-=2d+2N_-.
	\]
	第一式给出下界，第二式说明总数为偶数。\(d=0\) 时下界允许没有零点，也允许某个截面有符号相消的零点对。
\end{solution}

\begin{exercise}\label{g2:ex:disk-relative-euler-index}
	给闭圆盘 \(D^2\subset\mathbb C\) 取标准定向，在平凡定向二维丛上指定一个非零边界截面 \(v:S^1\to\mathbb C^*\)，令 \(k=\deg(v/|v|)\)。证明 \(v\) 可以延拓为处处非零的连续截面当且仅当 \(k=0\)。若一个光滑延拓 \(F:D^2\to\mathbb C\) 的零点全部孤立，证明其指标和为 \(k\)。当 \(v(z)=z^k\) 时，对任意整数 \(k\) 写出达到这个指标和的具体光滑延拓，并解释为何这不与平凡丛的绝对 Euler 类为零矛盾。
\end{exercise}
\begin{solution}
	若有非零延拓，将其除以长度便得到 \(D^2\to S^1\) 的延拓，故边界映射零伦，次数为零。反之，\(k=0\) 时 \(v/|v|:S^1\to S^1\) 零伦。用这个零伦沿半径填入圆盘，可得连续映射 \(g:D^2\to S^1\)，其边界值为 \(v/|v|\)。边界上的正函数 \(|v|\) 也有正延拓，例如取
	\[
		h(re^{\mathrm i\theta})=(1-r)+r|v(e^{\mathrm i\theta})|,
		\qquad h(0)=1.
	\]
	于是 \(hg\) 是所求非零延拓。

	设光滑延拓的零点为 \(z_1,\ldots,z_\ell\)，在每个零点周围取小闭圆盘 \(D_i\)，互不相交且不碰外边界。在挖去这些小圆盘内部的区域上，\(F/|F|\) 是到 \(S^1\) 的连续映射。该区域正向边界的同调类为外边界减去各小圆盘的正向边界；边界总类在区域的一维同调中为零。映入 \(H_1(S^1;\mathbb Z)\) 后便得
	\[
		k-\sum_{i=1}^\ell\deg\bigl((F/|F|)|_{\partial D_i}\bigr)=0.
	\]
	每个内边界次数按小圆盘本身的正向计算，正是对应零点的指标，故总和为 \(k\)。

	对 \(v(z)=z^k\)，若 \(k>0\)，取 \(F(z)=z^k\)；若 \(k<0\)，取 \(F(z)=\overline z^{\,-k}\)；若 \(k=0\)，取 \(F(z)=1\)。前两种延拓只有原点一个零点，其小圆周映射次数分别为 \(k\)，后一种没有零点。绝对 Euler 类属于 \(H^2(D^2;\mathbb Z)=0\)，只描述是否能自由选择一个非零截面；本题要求保持指定边界值，障碍属于 \(H^2(D^2,S^1;\mathbb Z)\cong\mathbb Z\)，恰由边界绕数 \(k\) 给出。
\end{solution}

\begin{exercise}\label{g2:ex:gysin-product-normalization}
	设 \(B\) 满足本章 Thom 同构的条件，\(r\ge2\)，并取平凡定向丛 \(E=B\times\mathbb R^r\)。令 \(u\) 是球面 \(S^{r-1}\) 的正向上同调生成元，仍用 \(u\) 表示它在 \(B\times S^{r-1}\) 上的拉回。证明
	\[
		H^k(B\times S^{r-1};R)=\pi^*H^k(B;R)\oplus\bigl(\pi^*H^{k-r+1}(B;R)\smile u\bigr),
	\]
	并计算两部分的沿纤维积分。写出上同调环的乘法约定，解释 Gysin 序列在这个例子中如何分裂。
\end{exercise}
\begin{solution}
	球面只有零次及 \(r-1\) 次上同调，且都是自由秩一 \(R\) 模。第4章的乘积比较因此给出所述直和；也可在球面的一份零胞腔、一份 \(r-1\) 胞腔的链模型上直接得到同一分解。正向生成元满足 \(u^2=0\)，因为它来自球面，而 \(H^{2r-2}(S^{r-1};R)=0\)。若 \(a\in H^\ell(B;R)\)，乘积中还有
	\[
		u\smile\pi^*a=(-1)^{(r-1)\ell}\pi^*a\smile u.
	\]
	这两条关系连同基底的乘法，给出分次张量积环 \(H^*(B;R)\otimes H^*(S^{r-1};R)\)。

	由正合性，\(\pi_!\pi^*a=0\)；由纤维归一化及投影公式，\(\pi_!(\pi^*b\smile u)=b\)。故
	\[
		\pi_!(\pi^*a+\pi^*b\smile u)=b.
	\]
	平凡丛有非零截面，所以 Euler 类为零，Gysin 序列分解成短正合列
	\[
		0\longrightarrow H^k(B;R)\xrightarrow{\pi^*}H^k(B\times S^{r-1};R)\xrightarrow{\pi_!}H^{k-r+1}(B;R)\longrightarrow0.
	\]
	映射 \(b\mapsto\pi^*b\smile u\) 是右逆，给出所述分裂。这个例子也检验了符号：基底类写在 \(u\) 前面时，积分结果就是 \(b\)。
\end{solution}

\begin{exercise}\label{g2:ex:gysin-sphere-product-base}
	在 \(B=S^2\times S^2\) 上，取两因子的正向生成元 \(x,y\)，于是
	\[
		H^*(B;\mathbb Z)=\mathbb Z[x,y]/(x^2,y^2),\qquad\deg x=\deg y=2.
	\]
	设定向实二平面丛 \(E\to B\) 满足 \(e_{\mathrm{Th}}(E)=ax+by\)，其中 \((a,b)\ne(0,0)\)，令 \(d=\gcd(|a|,|b|)>0\)。计算 \(H^*(S(E);\mathbb Z)\) 的各次群，并用 \(X=\pi^*x\)、\(Y=\pi^*y\) 以及适当的三次、五次生成元描述乘法。
\end{exercise}
\begin{solution}
	所有非零的 Euler 乘法由下面两个整数矩阵给出：
	\[
		H^0(B)=\mathbb Z\longrightarrow H^2(B)=\mathbb Z^2,\quad 1\longmapsto(a,b),
	\]
	\[
		H^2(B)=\mathbb Z^2\longrightarrow H^4(B)=\mathbb Z,\quad (s,t)\longmapsto bs+at.
	\]
	第一个映射单射，其余核为 \(\mathbb Z\oplus\mathbb Z/d\)，可由 Bézout 等式把列向量 \((a,b)\) 通过整数可逆行变换化成 \((d,0)\) 得到。第二个映射的像为 \(d\mathbb Z\)，核由 \((a/d,-b/d)\) 生成。基底的奇数次上同调为零，故 Gysin 序列在各次直接给出
	\[
		H^j(S(E);\mathbb Z)=
		\begin{cases}
			\mathbb Z,&j=0,3,5,\\
			\mathbb Z\oplus\mathbb Z/d,&j=2,\\
			\mathbb Z/d,&j=4,\\
			0,&\text{其余次数}.
		\end{cases}
	\]
	这里 \(\mathbb Z/1=0\)。二次群由 \(X,Y\) 生成，唯一加法关系为 \(aX+bY=0\)。四次群由 \(XY=\pi^*(xy)\) 生成，阶恰为 \(d\)，并有 \(X^2=Y^2=0\)。

	积分把三次群同构到上述第二个矩阵的核，把五次群同构到 \(H^4(B)\)。因此可选择唯一的生成元 \(z,w\)，使
	\[
		\pi_!z=\frac adx-\frac bdy,\qquad \pi_!w=xy.
	\]
	投影公式给出
	\[
		\pi_!(Xz)=x\smile\left(\frac adx-\frac bdy\right)=-\frac bdxy,\qquad
		\pi_!(Yz)=\frac adxy.
	\]
	由于五次积分是同构，得到 \(Xz=-(b/d)w\)、\(Yz=(a/d)w\)。其余正次数乘积由这些关系、分次交换律及五次以上群为零决定。这也展示了为什么只列出加法群还不够：同样的次数中，乘法仍保留 Euler 类两个分量的具体数值。
\end{solution}

\begin{exercise}\label{g2:ex:gysin-fiber-degree}
	设 \(E,E'\to B\) 为同秩 \(r\ge2\) 的定向实向量丛，\(B\) 连通且满足 Thom 同构的条件。设 \(F:S(E)\to S(E')\) 覆盖 \(\operatorname{id}_B\)，并在每条正向球面纤维上具有同一个整数映射度 \(m\)。证明
	\[
		e_{\mathrm{Th}}(E')=m\,e_{\mathrm{Th}}(E),\qquad
		\pi_!F^*b=m\,\pi'_!b.
	\]
	特别说明：纤维反向的同构会怎样改变 Euler 类和沿纤维积分。
\end{exercise}
\begin{solution}
	对 \(0<t\le1\)、\(v\in S(E)\)，定义 \(\overline F(tv)=tF(v)\)，并把各零向量送到相应的零向量。这是从 \((D(E),S(E))\) 到 \((D(E'),S(E'))\) 的连续映射，且覆盖基底恒等。零截面附近的连续性来自 \(\|\overline F(tv)\|=t\)；球面方向无论怎样变化，像的长度都趋于零。

	每条盘纤维上的相对映射度等于边界球面上的映射度 \(m\)，这是盘球对长正合列的自然性。故 \(\overline F^*u_{E'}\) 限制到每条纤维都是 \(m\) 倍正向生成元。由 Thom 同构在零次的唯一性，
	\[
		\overline F^*u_{E'}=m\,u_E.
	\]
	沿零截面拉回，得到 \(e_{\mathrm{Th}}(E')=m\,e_{\mathrm{Th}}(E)\)。对任意基底类 \(c\)，又有
	\[
		\overline F^*\Phi_{E'}(c)=p^*c\smile\overline F^*u_{E'}=m\Phi_E(c).
	\]
	空间对连接映射的自然性给出 \(\delta_EF^*b=\overline F^*\delta_{E'}b\)。两丛秩相同，积分定义中的次数符号也相同，因此代入上述等式便有 \(\pi_!F^*b=m\pi'_!b\)。

	纤维反向同构对应 \(m=-1\)，所以两边按所选方向计算的 Euler 类互为负数，积分也相差负号。若只反转同一丛的定向而保持底层映射为恒等，这正是 Thom 类、Euler 类和纤维积分同时反号的规则。
\end{solution}

\begin{exercise}\label{g2:ex:angular-form-thom}
	设 \(E\to B\) 为光滑定向实二平面丛，\(\pi:S(E)\to B\) 为其单位圆丛。给定一形式 \(\theta\in\Omega^1(S(E))\) 及闭二形式 \(\omega\in\Omega^2(B)\)，满足
	\[
		\int_{S(E)_x}\theta=1,\qquad \mathrm d\theta=\pi^*\omega.
	\]
	在 \(D(E)\setminus s_0(B)\) 上记径向投影为 \(\rho(v)=v/\|v\|\)。取光滑径向函数 \(\chi\)，使其在零截面附近为零，在球丛附近为一。证明
	\[
		U=\mathrm d(\chi\rho^*\theta)-p^*\omega
	\]
	给出实 Thom 类的相对形式代表，并求 \(e_{\mathrm{Th}}(E)\) 的实系数像。最后证明：若这个实 Euler 类非零，则球丛上不存在积分为一的闭一形式。
\end{exercise}
\begin{solution}
	因为 \(\chi\) 在零截面邻域恒为零，\(\chi\rho^*\theta\) 可通过补零光滑延拓到整个盘丛。闭性 \(\mathrm dU=0\) 来自 \(\mathrm d^2=0\) 及 \(\mathrm d\omega=0\)。在球丛上，\(\chi=1\)，所以
	\[
		i^*U=\mathrm d\theta-\pi^*\omega=0.
	\]
	因此 \(U\) 是相对闭形式。限制到单条盘纤维时，\(p^*\omega\) 为零；由外法向优先的 Stokes 公式，
	\[
		\int_{D(E)_x}U=\int_{S(E)_x}\theta=1.
	\]
	每条纤维上都是正向相对生成元，故 Thom 类的唯一性给出 \([U]=u_E\) 的实系数像。

	在零截面附近，\(U=-p^*\omega\)，于是
	\[
		e_{\mathrm{Th}}(E)_{\mathbb R}=[s_0^*U]=-[\omega].
	\]
	这个负号由相对代表在边界必须为零决定，与本节 Gysin 的定向约定一致。若另有闭一形式 \(\eta\) 满足每条纤维积分为一，则实系数 Gysin 序列给出 \(\pi_![\eta]=1\)，从而 \(e_{\mathrm{Th}}(E)_{\mathbb R}=0\)。这与假设矛盾，故所述闭形式不存在。
\end{solution}
